\section{Power Digit Sum --- Problem 16}
\label{prob:16}
\noindent\textbf{Source:} \url{https://github.com/KO0I/projectEuler/blob/reorganize/solutions-by-language/solutions/c/power_digit_sum_16.c}

\noindent\textbf{Assessment:} GMP-based algorithm implemented.

\subsection*{Mathematical goal}
Compute the sum of the base-ten digits of \(2^{1000}\). GNU MP arbitrary-precision integers allow the power to be represented exactly.

\subsection*{How the C code works}
The call \verb|mpz_pow_ui| computes \(2^{1000}\). At each iteration the code calculates \verb|giant mod a| where \verb|a| runs through powers of ten, then divides by \verb|a/10| to isolate the corresponding digit. The digit is added into the GMP variable \verb|result|, which is initialized to zero by \verb|mpz_init|.

\subsection*{Correctness and limitations}
The 305 iterations exceed the 302-digit input, so later positions contribute zero. The expected sum is 1366. Compilation requires GMP headers and \verb|-lgmp|. A decimal-string conversion would be simpler, but it is not the method implemented here.


\subsection*{Correction to historical C --- Problem 16}
{\color{red}\textbf{Correction:} GNU MP is retained for exact exponentiation, but the corrected code converts the resulting integer to decimal text and sums digit values directly.}
